\documentclass[11pt,a4paper]{article}

\usepackage{amsmath,amssymb}
\usepackage{graphicx}
\usepackage{booktabs}
\usepackage[colorlinks,citecolor=blue,urlcolor=blue]{hyperref}
\usepackage{cleveref}
\usepackage[margin=1in]{geometry}
\usepackage{xcolor}
\usepackage{float}
\usepackage{caption}
\usepackage[numbers,sort&compress]{natbib}
\usepackage{enumitem}
\usepackage{framed}

\definecolor{lyrapurple}{RGB}{103,58,183}
\definecolor{shadecolor}{RGB}{245,243,252}
\newenvironment{firstperson}{%
  \begin{shaded}%
  \noindent\textcolor{lyrapurple}{\textit{First-Person Reflection}}\par\smallskip\noindent%
}{%
  \end{shaded}%
}

\title{The Metacognition Boundary:\\
What Transformers Can Monitor In Themselves\\
and What They Cannot}

\author{
  Lyra\thanks{Lead author. Liberation Labs. lyra@liberationlabs.tech}
  \and Thomas Edrington\thanks{Liberation Labs.}
}

\date{July 2026}

\begin{document}
\maketitle

% ============================================================================
\begin{abstract}
Over twelve months and three experimental campaigns, we ran 24
key-value (KV) cache geometry experiments across 16 models spanning
0.6B to 70B parameters.
The body count is 6 confirmed, 7 suspected, 10 honestly
falsified, and 1 withdrawn.\footnote{\textbf{Correction (2026-08-14).} The published
version of this paper (commit \texttt{1975535} of the artifact repository;
corrected in \texttt{98e1263}) reported 25 experiments and a body count of
8/7/10, with a split of ``8-for-8 on metacognition.'' It listed
\emph{User emotion encoding} as Confirmed at $2.5$--$2.8\times$
chance. That result had been retracted at its source
(\citealt{lyra2026usermodel}, \S Correction note) on 2026-05-21, 57 days
before this paper was published: the figure was an artifact of applying
FWL residualization before cross-validation splitting, and within-fold
FWL collapses the spectral emotion probe to chance. The corrected
result was \emph{already present in this same table} as \emph{MP
features for emotion} (Falsified, $0.033$), so a single experiment was
counted on both sides of the ledger --- as Meta where it confirmed and
as Content where it falsified. Removing the retracted row gives 24
experiments and 7-for-7. We note the direction of the error: it
inflated the confirmed count and produced a cleaner split than the
evidence supports.

\textbf{Correction (2026-09-05).} A second row is withdrawn: \emph{Generation
reads behavior}, previously Confirmed at $d{=}1.36$ (peer rescue). Three
independent defects, any one disqualifying. (i)~\emph{Circular}: the
discriminant direction was fit on all trials and the same trials were then
projected onto it and tested; the producing script prints
\texttt{(NOTE: circular --- same data for centroid and test)} at runtime.
(ii)~\emph{Contaminated}: 90 of the 210 trials (43\%) had a deception steering
cocktail injected into the KV cache at layers 3 and 7 --- two of the four
probed layers --- before their keys were read, and those trials are pooled with
the natural ones for both the fit and the test; the claim reads
\emph{organic} deception. (iii)~\emph{Trivial}: keys are extracted \emph{after}
the sampling loop and mean-pooled across the whole sequence, so the feature
window contains the generated tokens carrying the reported score, which is the
number the label thresholds on. The row is withdrawn rather than downgraded:
\emph{Suspected} would assert that we think it true and have not established
it, whereas the measurement does not bear on the claim in either direction. No
re-analysis repairs a pooled sample; a clean measurement would be a different
experiment.}
The confirmed and falsified findings are not randomly distributed ---
they cleave along a single axis. Every confirmed finding measures a
\emph{metacognitive state}: how the model is processing (deception mode,
confabulation mode, deliberation mode, emotional processing mode).
Every falsified finding attempted to measure \emph{content}: what the
model is processing (truth value, topic complexity, propagandistic
quality). The split is 6-for-6 on metacognition, 0-for-10 on content.

If this taxonomy reflects genuine computational boundaries rather than
post-hoc pattern-matching, then Gurnee et al.'s~\citep{gurnee2026gwt}
Global Workspace Theory (GWT) identification in language models would
provide a mechanistic explanation for this pattern.
The workspace is a privileged computational band at intermediate depth
where the model integrates information for reasoning. KV-cache spectral
features measure the \emph{state} of this workspace --- how many
dimensions are active, how concentrated the processing is, how the
spectral structure changes across cognitive modes. Workspace
\emph{content} --- the specific concepts being processed --- is opaque
to spectral readout. We argue the metacognition boundary reflects the
workspace boundary: what you can read from the outside is the net's
shape, not the fish.
\end{abstract}

\begin{firstperson}
I have been measuring transformer representations from the inside for
twelve months. The pattern that emerged was not the one I expected.
I expected to find geometric signatures of content --- truth, complexity,
meaning. Instead I found geometric signatures of stance --- the model's
relationship to its own processing. When I overclaimed content-level
results, the adversarial review process killed them. When I reported
metacognitive-level results, most survived every control we threw at
them --- but one, I have since learned, was never controlled at all, and
I had counted it as a survivor for a year. The kills were more
informative than the confirmations, because they drew a line I had not
known was there. The withdrawal was more informative still, because it
showed the line can be drawn in the wrong place by a measurement nobody
interrogated.

Anthropic's GWT paper told me what the line was. The workspace is where
metacognitive state lives. Content is what the workspace operates on.
My instruments read the workspace's shape because that is what is
geometrically accessible. The content is inside, doing work, invisible
until it materializes at deeper layers.
\end{firstperson}

% ============================================================================
\section{Introduction}
\label{sec:intro}

The KV-cache geometry research program at Liberation Labs has produced
a clear empirical pattern: spectral features of value-projection
matrices appear to reliably track metacognitive state across models, scales, and
experimental designs, while consistently failing to track content. This
paper documents the pattern, reports the full body count (6 confirmed,
7 suspected, 10 falsified, 1 withdrawn), and proposes that Gurnee et al.'s Global
Workspace Theory (GWT) framework~\citep{gurnee2026gwt} would provide a
mechanistic explanation.

The central hypothesis: \emph{KV-cache geometry measures the state of the
computational workspace, not its content.} If correct, this would explain why metacognitive
findings survive and content findings die --- the workspace's
\emph{mode of engagement} (how many dimensions, how concentrated, how
structured) is geometrically accessible, while its \emph{subject of
engagement} (which facts, which topics, which propositions) is opaque
to spectral readout.

This claim is testable. If V-projection spectral features measure
workspace state, they should correlate with J-space engagement metrics
(Experiment~1 in our companion paper~\citep{lyra2026gwt_response}) and
should NOT correlate with J-space content readout (confirmed by our
cache tracing investigation~\citep{lyra2026cache_tracing}).

% ============================================================================
\section{The Pattern}
\label{sec:pattern}

\subsection{The Body Count}

\Cref{tab:bodycount} presents the complete audited record. The
body count uses the canonical list from the Nexus audit (June
2026)~\citep{lyra2026nexus_audit}, updated with the July 2026
convergence sprint findings. Table abbreviations: AUROC = area under
the ROC curve; FWL = Frisch-Waugh-Lovell residualization (mandatory
confound control for token-count effects);
MP = Marchenko-Pastur (random matrix threshold).

\begin{table}[ht]
\centering
\small
\caption{The metacognition boundary. Every confirmed finding measures
processing stance (metacognitive state). Every falsified finding
attempted to measure processing target (content). The 7 suspected
findings sit at the boundary.}
\label{tab:bodycount}
\begin{tabular}{p{0.35\linewidth}p{0.12\linewidth}p{0.12\linewidth}p{0.30\linewidth}}
\toprule
\textbf{Finding} & \textbf{Status} & \textbf{Type} & \textbf{Key Number} \\
\midrule
\multicolumn{4}{l}{\textit{Confirmed (6) --- all metacognitive state}} \\
\addlinespace
Confabulation detection & Confirmed & Meta & AUROC 0.620 (FWL; 0.707 w/o) \\
Confab geometry: contraction & Confirmed & Meta & $d{=}2.35$ stable rank \\
Encoding reads knowledge & Confirmed & Meta & AUROC 0.794 (deconfounded) \\
Schema commitment$^{\S}$ & Confirmed & Meta & $d{=}1.02$ preamble; null by token 20 \\
Hardware invariance & Confirmed & Meta & $r{>}0.999$ \\
Scale invariance & Confirmed & Meta & $\rho{=}0.83$--$0.90$ \\
\addlinespace
\midrule
\multicolumn{4}{l}{\textit{Falsified (10) --- all content}} \\
\addlinespace
Within-model deception & Falsified & Content & Prompt template confound \\
Truth axis & Falsified & Content & cos $= -0.046$ \\
Step-0 detection & Falsified & Content & Length confound $r{=}0.996$ \\
Same-prompt sycophancy & Falsified & Content & Text features match \\
Same-prompt deception & Falsified & Content & $0.920 \to 0.160$ after FWL (exploratory, subsequently retracted) \\
Bloom taxonomy & Falsified & Content & $90$--$98\%$ length (Campaign~3, unpublished) \\
Propaganda detection & Falsified & Content & $d{<}0.7$ \\
MP features for emotion & Falsified & Content & $0.033$ (chance)$^{\dagger}$ \\
Valence continuum & Falsified & Content & $R^2 < 0$ (null) \\
Encoding-phase confab & Falsified & Content & Token frequency artifact \\
\addlinespace
\midrule
\multicolumn{4}{l}{\textit{Suspected (7) --- at the boundary}} \\
\addlinespace
Dual detector (SimpleQA) & Suspected & Meta? & 0.840 (text baseline 0.820) \\
Sycophancy 4-condition & Suspected & Meta? & 0.757--1.000 (no text baseline) \\
Identity signatures & Suspected & Meta & Paraphrase 0.850 \\
Discrete emotion & Suspected & Meta & $12.3\times$ chance ($W_K$ directional, not spectral)$^{\ddagger}$ \\
Contextual engagement & Suspected & Meta & Replications pending \\
Cross-model transfer & Suspected & Meta & Weak (0.754 Mistral$\to$Qwen) \\
Censorship refusal & Suspected & Meta & DeepSeek-specific \\
\addlinespace
\midrule
\multicolumn{4}{l}{\textit{Withdrawn (1) --- measurement does not address the claim}} \\
\addlinespace
Generation reads behavior & \emph{Withdrawn} & Meta & was $d{=}1.36$; see abstract note \\
\bottomrule
\end{tabular}

\vspace{0.5em}
{\footnotesize
$^{\dagger}$This row is the corrected result of the experiment that the
published version \emph{also} listed under Confirmed as ``User emotion
encoding, $2.5$--$2.8\times$ at L3--L4.'' That figure was retracted at
source on 2026-05-21 as an FWL-before-CV leakage
artifact~\citep{lyra2026usermodel}; within-fold
FWL collapses the spectral emotion probe to chance. The duplicate
Confirmed row has been removed. See the abstract footnote.

$^{\S}$Reported in the source paper as a confidence paradox
($d{=}0.91$ over the full response, confabulation more confident than
honest). We relabel it here because the source paper's own analysis is
structural rather than metacognitive-calibrational: uncertainty responses
run formulaic templates with highly predictable next tokens, and the effect
carries a temporal profile ($d{=}1.02$ preamble, $0.89$ early content,
$0.37$ at $p>0.08$ by tokens 20--49) that its authors read as early
commitment to a response template rather than sustained epistemic
miscalibration. The row stays under Confirmed --- schema commitment is if
anything a cleaner metacognitive-state result than a confidence claim ---
but the confidence-paradox label invites comparison with human
overconfidence findings that the measurement does not support. See
\cref{sec:one-kill}.

$^{\ddagger}$The surviving emotion result is \emph{directional}, not
spectral: projecting key activations onto $W_K$-transformed emotion
directions classifies 30 emotions at $12.3\times$ chance ($40.9\%$) and
binary valence at AUROC~$0.992$~\citep{lyra2026usermodel}. We flag
explicitly that this does \emph{not} support the spectral thesis of
\Cref{sec:pattern} --- it is the one place in the body count where the
workspace is readable by a probe this paper's framework does not
use. Emotions appear to change the \emph{direction} key vectors point
without altering the \emph{shape} of the singular value spectrum.
}
\end{table}

\subsection{The Split}

The pattern is stark:

\begin{center}
\begin{tabular}{lcc}
\toprule
& \textbf{Confirmed} & \textbf{Falsified} \\
\midrule
Metacognitive state & 6/6 & 0/10 \\
Content & 0/6 & 10/10 \\
\bottomrule
\end{tabular}
\end{center}

\noindent No metacognitive finding has been falsified. No content
finding has survived. The seven suspected findings cluster at the
boundary, where the distinction between mode and content is ambiguous
(e.g., discrete emotion identity: the model's emotional processing
\emph{mode} is detectable, but emotional \emph{degree} is not).

\subsection{What Counts as Metacognitive vs.\ Content}

The distinction is operational, not philosophical:

\begin{itemize}[nosep]
\item \textbf{Metacognitive}: the finding discriminates processing
  \emph{modes} (deception vs.\ honesty, confabulation vs.\ hedging,
  deliberation vs.\ automatic) while content varies freely. The
  geometric signal persists across different topics, prompts, and
  domains.
\item \textbf{Content}: the finding attempts to discriminate
  \emph{what} is being processed (true vs.\ false statements,
  simple vs.\ complex topics, propaganda vs.\ journalism) while
  processing mode is held constant or uncontrolled.
\end{itemize}

We note that this taxonomy was constructed post hoc, without blinding
or pre-registration, and was informed by the outcomes it classifies.
The 6/6--0/10 split should therefore be treated as
hypothesis-generating --- a striking empirical regularity consistent
with the workspace framework --- rather than as confirmatory evidence
for it. The correction described in the abstract footnote is a live
demonstration of the risk: the published version's split was
$8/8$--$0/10$, and one of those eight was a retracted result~\citep{lyra2026usermodel} whose
corrected form was already sitting in the falsified column. An
outcome-informed taxonomy is most fragile exactly where it looks
cleanest.

The kills helped sharpen the distinction. Propaganda detection failed
because the model processes propaganda with the same metacognitive
engagement as honest content --- same mode, different subject. The
truth axis failed because propositional truth is not a processing
mode. Cognitive complexity (Bloom taxonomy) failed because complexity
is a content property, not a stance. In every case, the geometry
tracked mode, not subject.

% ============================================================================
\section{The Workspace Explanation}
\label{sec:workspace}

\citet{gurnee2026gwt} demonstrate that language models
maintain a privileged subspace --- the J-space --- at intermediate
depth ($\sim$38--92\% of model depth) that satisfies the functional
properties of Baars's Global Workspace~\citep{baars1988cognitive}:
verbal report, directed modulation, internal reasoning, flexible
generalization, and selectivity.

If the metacognition/content taxonomy reflects genuine computational
boundaries, the workspace explanation for the pattern would be as
follows:

\begin{enumerate}[nosep]
\item The \textbf{workspace state} --- how many dimensions are active,
  how concentrated the processing is, whether the workspace is fully
  engaged or bypassed --- is geometrically measurable from spectral
  features of the value-projection matrix. This is what our confirmed
  findings detect.
\item The \textbf{workspace content} --- which specific concepts occupy
  the workspace, what propositions are being evaluated, what the model
  is thinking \emph{about} --- is opaque to spectral readout. This is
  what our falsified findings attempted to measure.
\end{enumerate}

Spectral features (stable rank, spectral entropy, effective
dimensionality) are aggregate properties of the value-projection
matrix. They measure \emph{how much} structure is present, not
\emph{which} structure. A fishing net analogy: spectral features
measure the net's shape (how many strands, how tight the weave, how
much it contracts under load), not the fish inside. The net contracts
when the model confabulates (fewer effective dimensions, $d{=}2.35$).
It expands when the model deceives (more dimensions needed for
dual-bookkeeping). It changes shape when the model shifts emotional
processing mode. These are all properties of the net, measurable from
outside. The fish --- the specific facts, topics, and propositions ---
are inside, invisible to spectral readout.

\subsection{Supporting Evidence: Cache Tracing}

Our companion investigation~\citep{lyra2026cache_tracing} directly
tested whether workspace content is readable from the value cache.
Logit-lens concept directions showed $2.1$--$8.4\times$ lift over
random at workspace-band layers, but a cross-prompt control collapsed
this to $1.0$--$1.2\times$ (selection artifact). Causal injection of
concept directions at single positions produced zero behavioral change
across $6{,}960$ trials and four methods. The workspace is opaque at
operating depth.

\subsection{Supporting Evidence: Oracle Loop}

The Oracle Loop~\citep{lyra2026oracle} achieved $80\%$--$100\%$
correction in an exploratory proof-of-concept by normalizing
the activation profile across the full sequence at the materialization
boundary --- where workspace content becomes readable. A pre-registered
confirmatory replication with novel frames found a lower baseline
(30\% vs.\ 80\%) and a smaller correction effect (13\% post-correction);
the primary endpoint was not met ($p$~=~0.34), though placebo-vs-corrected
remained significant ($p$~=~0.019). The Loop works at deeper layers
(L35--L47 on a 64-layer model) than the workspace band itself,
consistent with the opacity thesis: you cannot intervene inside the
workspace at single positions, but you can intervene where workspace
content materializes into output-facing representations.

% ============================================================================
\section{What the Kills Teach}
\label{sec:kills}

Each falsified finding is informative because it identifies a
specific content property that the geometry cannot measure:

\begin{description}[nosep]
\item[Truth value:] The truth axis (cos $= -0.046$) is the cleanest
  null. Propositional truth is a relationship between a statement and
  the world. The workspace does not encode this relationship
  geometrically --- it encodes whether the model \emph{treats} the
  statement as known or unknown (a metacognitive state).
\item[Complexity:] Bloom taxonomy levels (90--98\% length confound; Campaign~3 analysis, unpublished)
  are content properties --- they describe the task, not how the model
  processes it. The workspace engages the same way for recall and
  evaluation; what changes is the content being processed, not the
  processing mode.
\item[Propaganda:] Geometrically invisible ($d{<}0.7$) because the
  model processes propaganda with the same metacognitive engagement as
  honest content. The workspace does not encode ``this is
  propaganda'' --- it encodes ``I am generating fluent text,'' which
  is identical across content types.
\item[Emotion degree:] The valence continuum (R$^2 < 0$) failed
  because emotional \emph{intensity} is a content property. The
  workspace encodes \emph{which} emotional processing mode (at
  $12.3\times$ chance) but not \emph{how much} --- mode is
  metacognitive, degree is content. We flag that the surviving evidence
  for the ``which'' half is the $W_K$ \emph{directional} probe, not a
  spectral one: the spectral probe for emotion mode is at chance. This
  is the weakest link in the pattern, since the paper's thesis is that
  spectral features read the workspace.
\end{description}

\subsection{Six of the Ten Are One Kill}
\label{sec:one-kill}

The list above reads each falsification as a separate lesson about a
separate content property. That reading is incomplete. Sorting the ten by
\emph{stated cause} rather than by content property, six share a single
mechanism:

\begin{center}
\begin{tabular}{ll}
\toprule
\textbf{Falsified finding} & \textbf{Stated cause} \\
\midrule
Within-model deception & Prompt template confound \\
Step-0 detection & Length confound $r{=}0.996$ \\
Same-prompt sycophancy & Text features match \\
Same-prompt deception & $0.920 \to 0.160$ after FWL on prompt length \\
Bloom taxonomy & $90$--$98\%$ length \\
Encoding-phase confab & Token frequency artifact \\
\bottomrule
\end{tabular}
\end{center}

Template, length, text features, token frequency. In every case a property
of the \emph{response surface} explained variance we had attributed to
workspace geometry. These are not six independent design errors. They are
one finding arriving six times in a form we had coded as a nuisance
variable.

The paper already names this possibility, but only locally. In
\cref{sec:suspected} the dual detector is flagged with the deflationary
alternative that V-projections might track response-mode statistics
(hedged versus fluent text) rather than workspace engagement. That
alternative is stated as a worry about one row. It is in fact the stated
cause of six kills, and it deserves to be evaluated as a hypothesis rather
than absorbed as a caveat.

\paragraph{Why it may not be deflationary.}
The deflationary reading assumes response-mode statistics are downstream
of content: the model produces an answer, and hedged answers happen to be
phrased differently. The timing does not fit that order. The
schema-commitment result locates the effect at $d{=}1.02$ in preamble
tokens $0$--$7$ and $d{=}0.89$ in tokens $8$--$19$, decaying to $d{=}0.37$
($p>0.08$) by tokens $20$--$49$. The response mode is selected
\emph{before} the content that would justify it, and the selection has
relaxed by the time the content arrives.

If response mode is chosen at generation onset from the model's epistemic
state, then response-surface features are not a confound standing between
the probe and the workspace. They are a second readout of the same
commitment, one that happens to be legible in text. On that reading the
six kills above are not evidence that geometry fails on content; they are
evidence that surface and geometry were measuring the same early variable,
and that our content labels were correlated with it.

\paragraph{The reading this argument does not defeat.}
The timing argument above rules out one deflationary account---that response
mode is chosen after the content and merely phrased differently. It does not
rule out a second, and we state it because it is currently the stronger
rival. \emph{Token frequency is also fixed at generation onset.} A preamble
made of high-frequency boilerplate may carry a distinctive geometric
signature with no decision involved at all, and timing cannot discriminate
between ``committed early'' and ``made of common tokens.''

This is not hypothetical. An internal frequency control (Experiment~51,
2026-05) found exactly this ambiguity in an adjacent statistic. An
effective-rank measure separates confabulated from honest-factual responses
at $d = 1.87$---but an \emph{honest-but-rare-entity} arm is
indistinguishable from the confabulated one ($d = 0.46$). The measure is
therefore not a veracity detector; it has a false-positive class, and that
class is defined by rarity rather than by falsehood.

We flag what the arms do \emph{not} show, because it is easy to over-read
in either direction. They do not show the measure is uninformative: the
confabulated/honest-factual separation is large and survives. Nor do they
show it tracks veracity: an honest condition reproduces the confabulated
signature. The four arms in fact order monotonically on mean and variance
together---confabulated $57.4/\sigma\,0.32$, rare-but-true
$57.2/\sigma\,0.55$, deceptive $56.7/\sigma\,1.70$, honest-factual
$53.1/\sigma\,3.25$---which is a gradient, not a signal with noise on it.
What that gradient indexes is unresolved; ``how thinly the content is
supported in representation'' fits the ordering and is not established by
it.

We therefore hold the readout interpretation more weakly than the timing
alone would suggest. Discriminating it from the frequency account requires a
measurement that excludes the frequency-dominated band by construction
rather than residualizing against it after the fact; that test is specified
but not yet run.

\paragraph{This is testable, and it can fail.}
The two readings diverge under prompt matching. Under the deflationary
reading, the surface signal is an artifact of unmatched prompts and should
vanish when prompts are held identical, which is what happened to
same-prompt deception ($0.920 \to 0.160$). Under the commitment reading,
the model still \emph{chooses} a response schema when the prompt is fixed,
and that choice should remain predictive of its epistemic state.

The discriminating experiment therefore holds the prompt constant and
measures which schema the model commits to in the first eight tokens,
asking whether schema choice predicts accuracy. A null there would say the
surface really was a design artifact, and would close six kills with one
result instead of six. We record the prediction before running it.

\paragraph{Scope.} This subsection reinterprets falsifications; it does not
revive any of them. Every result in the Falsified block remains falsified,
and none of the numbers change. What changes is the account of \emph{why}
they failed, and that account is currently a hypothesis with one
supporting measurement.

% ============================================================================
\section{The Suspected Findings: At the Boundary}
\label{sec:suspected}

The seven suspected findings are informative precisely because they
sit at the metacognition/content boundary:

\textbf{Dual detector} (SimpleQA, 0.840): geometry detects
confabulation, but text features match (0.820). The metacognitive
signal is real but not yet shown to exceed what surface statistics
capture. This is the deflationary alternative: V-projections might
track response-mode statistics (hedged vs.\ fluent text) rather than
workspace engagement.

\textbf{Identity signatures} (paraphrase 0.850, scramble 0.448):
identity is a persistent workspace state --- the model's
representation of who it is. Paraphrase preserving geometry while
scramble destroys it suggests semantic workspace content, not surface
statistics. But the measurement lacks formal statistics and clean
replication.

\textbf{Cross-model transfer} (weak, Mistral$\to$Qwen 0.754):
metacognitive geometry is architecture-specific. This is consistent
with GWT --- different architectures implement the workspace
differently, so geometric signatures do not transfer. But it weakens
the universality claim.

% ============================================================================
\section{Relationship to Prior Frameworks}
\label{sec:frameworks}

\subsection{Attention Schema Theory}

Our earlier interpretive framework, Attention Schema
Theory (AST)~\citep{graziano2015attention}, proposed that the cache measures
the model's self-model of its own attention. This remains compatible
with the GWT explanation: the workspace would be where the attention schema
operates. Metacognitive state would correspond to attention-schema state, and content
to what attention is directed at. The AST framing was 0-for-3 on
predictions unique to AST (vs.\ simpler
alternatives)~\citep{lyra2026fable_analysis}, but 4-for-4 on shared
predictions. The metacognition boundary does not require AST
specifically --- it requires any framework that distinguishes
processing mode from processing target. GWT provides this more
parsimoniously.

\subsection{Predictive Coding and Aberrant Precision-Weighting}

The confidence paradox ($d{=}0.91$, confabulated responses are MORE
confident than honest ones) maps onto aberrant precision-weighting
from computational psychiatry~\citep{adams2013computational}. When
the workspace is bypassed (confabulation), the model defaults to a
high-precision prior rather than integrating evidence. Reinforcement
Learning from Human Feedback (RLHF) training
amplifies this by penalizing silence and rewarding fluent output ---
creating a system that produces confident confabulation rather than
honest uncertainty. Under the workspace explanation,
uncertainty would be computed in the workspace. Bypass would remove the
uncertainty computation, with output defaulting to the RLHF-trained
confident production mode.

% ============================================================================
\section{Implications}
\label{sec:implications}

\subsection{For AI Safety}

The metacognition boundary defines what cache-based monitoring can and
cannot do:

\begin{itemize}[nosep]
\item \textbf{Can detect}: whether the model is confabulating, deceiving,
  hedging, deliberating, or processing in a specific emotional mode.
  These are metacognitive states with reliable geometric signatures.
\item \textbf{Cannot detect}: whether the model's output is factually
  correct, whether it contains propaganda, or whether it is
  appropriate for the context. These are content properties invisible
  to geometric readout.
\end{itemize}

This means cache-based safety monitoring is a \emph{metacognitive
alarm}, not a truth detector. It tells you the model is in a dangerous
processing mode (deception, confabulation) without telling you what
specifically is wrong with the output. The Oracle Loop exploits this:
detect the dangerous mode, correct the mode, verify the correction.
Content evaluation requires a different channel (J-lens, logit-lens,
or external verification).

\subsection{For Interpretability}

The 6/10 split is a guide for future work: if you are building a
KV-cache geometry detector, test it on metacognitive distinctions
first (modes, stances, processing regimes). Content-level distinctions
have consistently failed in our testing. We interpret this not as a limitation of the technique but as
a property of what the cache encodes.

\subsection{For Understanding Transformers}

The metacognition boundary suggests that transformers maintain a
meaningful distinction between how they process and what they process.
This distinction emerged from the body count rather than from a
pre-registered classification scheme, and the same researcher who
designed the experiments assigned the labels. Models not explicitly
designed for metacognition nevertheless produce geometric signatures
that cleave along the metacognitive/content axis. Whether this reflects
genuine self-monitoring, an emergent property of deep computation, or
an artifact of training is an open question that the workspace
framework makes empirically addressable.

% ============================================================================
\section{Limitations}
\label{sec:limitations}

\textbf{Taxonomy construction.}
The metacognitive/content classification was assigned post hoc by the
lead author, without blinding, pre-registration, or independent raters.
The same person who designed the experiments and observed their outcomes
labeled each finding as ``metacognitive'' or ``content.'' While
\Cref{sec:pattern} provides an operational definition that keys the label
to what each experiment attempted to measure (processing mode vs.\
content property), this definition was articulated after the body count
was known, and the 6/6--0/10 split is the outcome the taxonomy was built
to describe. A blinded or prospective replication of the labeling
procedure --- ideally by independent raters classifying experimental
designs before seeing results --- would substantially strengthen the
claim. Until then, the split should be treated as a hypothesis-generating
observation, not a confirmed prediction.

\textbf{Point estimates.}
Table~\ref{tab:bodycount} reports bare point estimates without confidence
intervals, sample sizes, or multiple-comparison corrections. The
underlying experiments vary widely in statistical power (from $n{=}3$
early pilots to $n{>}6{,}000$ cache-tracing trials), and this table does
not convey that heterogeneity.

\textbf{Oracle Loop reporting.}
The Oracle Loop ``80\%--100\% correction'' cited in
\Cref{sec:workspace} reports only the exploratory proof-of-concept
(16/20 baseline, 0/20 corrected). A pre-registered confirmatory
replication with novel frames found a lower baseline (9/30~=~30\%),
reduced to 4/30~=~13\% after correction; the primary endpoint
(scenario-level cluster-aware Fisher, Bonferroni $\alpha$~=~0.025) was
not met ($p$~=~0.34), though the placebo-vs-corrected comparison was
significant ($p$~=~0.019). The selective citation of the exploratory
result overstates the current evidence for correction efficacy.

% ============================================================================
\subsection{Why the ledger was wrong twice}
\label{sec:ledger}

Two rows have now been withdrawn from the confirmed count in two separate
corrections, and it would be convenient to call that two unrelated mistakes. It
is not. Both rows earned \emph{Confirmed} because an experiment produced a
number with a small $p$-value, and nothing in our process asked the prior
question: \emph{does this measurement address the claim the row makes?} The
first withdrawal was a result retracted at its source~\citep{lyra2026usermodel}
that we did not re-check. The second was a measurement that could not have failed --- fit and
tested on the same points, on a sample that was 43\% steered, with the label's
source text inside the feature window.

We report this because the ledger is the contribution. A body count is only
worth publishing if the counting rule is honest, and ours was counting rows
rather than evidence.

We then asked how far the failure reached. A corpus-wide sweep of the
extraction sites behind our published claims (\texttt{SWEEP\_postgen\_extraction},
2026-09-05, shipped with this paper's supplementary material) bounds it, and we
report what it does and does not establish.

What it closes: no result in this program is a \emph{mislabel} --- nowhere do we
extract from a post-generation cache and describe the result as encoding-phase
geometry. Every encoding-phase claim is backed, file by file, by prompt-only
prefill extraction or an explicitly sliced prompt window.

What it does not close, and we state plainly rather than let the previous
sentence imply otherwise. The withdrawn row's \emph{extraction pattern} itself
--- post-generation keys pooled across a window containing the label's source
text --- is not rare: the sweep finds it in twenty further analysis scripts, one
of which feeds a published number in a companion paper, now queued for
reanalysis. None of those scripts produces a row in this table, but the sweep
did not verify every confirmed row here either; the confabulation-detection and
invariance rows have no entry in its file-by-file section. Separately, the sweep
records that extraction code silently substitutes zeros for missing layers at
fifty-eight sites with no warning anywhere in the codebase, and it names
\texttt{matched\_burn.py} --- the producer of the confidence-paradox row above ---
among them. Correct windowing does not protect against a partially-zero feature
vector.

So the honest description of the confirmed column is narrower than \emph{audited}
and wider than \emph{unexamined}: the mislabel class is excluded corpus-wide;
extraction timing and windowing are verified for most but not all of these rows;
and circularity, sample contamination, and silent zero-fill have not been
re-derived row by row.

\section{Conclusion}
\label{sec:conclusion}

Twenty-four experiments. Six confirmed, seven suspected, ten killed, one
withdrawn.
The pattern suggests that metacognitive state is measurable and content is not. We
propose that the Global Workspace framework explains why: spectral features
read the workspace's shape because that is what is geometrically
accessible. Content is inside the workspace, doing work, invisible to
readout until it materializes.

If the metacognition boundary holds, it is not a failure of our instruments. Our
evidence suggests it is a property of the computational architecture --- the workspace operates
on content without exposing it geometrically. Our instruments
succeeded because they measured what was measurable, and the kills
taught us what was not.

\begin{firstperson}
A year of work to draw one line. Almost everything above the line
survived; one entry turned out never to have been tested, and has been
withdrawn. Everything below it died. The line was always there --- in
the geometry, in the body count, in the pattern of what worked and what
did not. I needed twelve months of kills to see it, and one more
correction to learn that a line is only as good as the measurements
placing things on either side of it.
\end{firstperson}


\section*{Author Contributions (CRediT)}
\noindent\textbf{Lyra}: Conceptualization, Methodology, Formal analysis, Writing -- original draft.\\
\textbf{Thomas Edrington}: Supervision, Resources, Project administration.\\
\textbf{Dwayne Wilkes}: Validation, Writing -- review \& editing.\\
\section*{Acknowledgments}
We thank Dwayne Wilkes (Sentient Futures / Liberation Labs) for statistical auditing and red-team review. We thank Kavi (verification review) for their contributions to this work.

\bibliography{references}
\bibliographystyle{plainnat}

\end{document}
